% ============================================================
%  sec-4-training-pipeline.tex  —  4.4 Training Pipeline
% ============================================================

\section{Fine-Tuning Pipeline Implementation}
\label{sec:training-pipeline}

\subsection{Patched Training Entry Point}
\label{sec:train-entry-point}

The fine-tuning entry point is \filepath{train\_sft\_ascend.py},
a modified version of the original \texttt{stage1\_sft/train.py}.
It preserves the complete original training logic while applying
six sequential patches at startup. The order of these patches is
critical: \texttt{torch\_npu} must be imported before any other
package, because it modifies PyTorch's device registry at import
time.

\tabref{patch-sequence} documents the six patches in execution order.

\begin{table}[htbp]
  \centering
  \caption{Ascend patch sequence in \texttt{train\_sft\_ascend.py}}
  \label{tab:patch-sequence}
  \begin{tabularx}{\textwidth}{lXX}
    \toprule
    \tableheader{Step & Action & Why} \\
    \midrule
    1 & \texttt{import torch\_npu;
        transfer\_to\_npu()}
      & Registers \texttt{"npu"} device; intercepts all
        \texttt{.cuda()} calls globally \\
    \rowcolor{TableRowOdd}
    2 & Patch \texttt{mamba\_ssm.ops.selective\_scan\_interface}
      & Replaces CUDA kernel before model is constructed \\
    3 & Patch \texttt{MoETokenDispatcher.forward}
      & Replaces CUDA MoE dispatch before model is constructed \\
    \rowcolor{TableRowOdd}
    4 & Call \texttt{\_patch\_autocast()}
      & Replaces \texttt{torch.cuda.amp.autocast} with
        \texttt{torch.npu.amp.autocast} \\
    5 & Import Megatron-Bridge and Nemotron modules
      & Safe to import after patches are in place \\
    \rowcolor{TableRowOdd}
    6 & Set \texttt{TORCH\_DISTRIBUTED\_DEFAULT\_BACKEND="hccl"}
      & Ensures \texttt{init\_process\_group()} uses HCCL \\
    \bottomrule
  \end{tabularx}
\end{table}

The patches register the Ascend device, replace the Mamba and MoE kernels
before model construction, and ensure the distributed backend and autocast
mechanisms are compatible with the NPU hardware.

\subsection{SFT Training Configuration}
\label{sec:sft-config}

The production configuration (\filepath{configs/sft\_ascend.yaml})
targets a single Ascend 910C node with 8~NPUs, fine-tuning for
3,400 iterations — approximately one epoch over the 137K-sample
dataset. Key parameters include a global batch size of 32, a learning rate
of $2.0 \times 10^{-5}$, and a tensor model parallelism size of 4 to
optimise memory usage across the NPUs.

\tabref{config-rationale} explains the rationale for key
configuration choices.

\begin{table}[htbp]
  \centering
  \caption{SFT configuration parameter rationale}
  \label{tab:config-rationale}
  \begin{tabularx}{\textwidth}{lX}
    \toprule
    \tableheader{Parameter & Rationale} \\
    \midrule
    \texttt{tensor\_model\_parallel\_size: 4}
      & Splits the 3,008-dimensional hidden weight matrices across 4
        NPUs, keeping per-NPU memory below 55~GB \\
    \rowcolor{TableRowOdd}
    \texttt{context\_parallel\_size: 2}
      & Splits the 4,096-token sequence across 2 NPU groups for
        attention computation, reducing peak activation memory \\
    \texttt{packed\_sequence: true}
      & Uses the packed Parquet shards from SFT data prep, ensuring
        no padding tokens are included in the loss \\
    \rowcolor{TableRowOdd}
    \texttt{calculate\_per\_token\_loss: true}
      & Required for packed sequences: normalises the cross-entropy
        loss by the number of non-masked tokens rather than by batch
        size \\
    \texttt{peft: null}
      & Disables LoRA/adapter fine-tuning; full weight update
        achieves better domain adaptation for the telecom use case \\
    \rowcolor{TableRowOdd}
    \texttt{finetune: true}
      & Skips loading the pre-training optimiser state from the
        checkpoint, allowing the fine-tuning run to begin with a
        fresh optimiser \\
    \bottomrule
  \end{tabularx}
\end{table}

\subsection{Test Configuration}
\label{sec:test-config}

The test configuration (\filepath{configs/sft\_tiny.yaml}) enables
rapid validation on a single \acrshort{npu} in approximately
15 minutes. It substitutes a smaller Qwen3-1.8B recipe with 100 training
iterations and a global batch size of 8, allowing for complete pipeline
validation — including data loading and W\&B logging — without the
parallelism complexity of the full model.

\subsection{W\&B Artifact Lineage}
\label{sec:loss-validation}

The fine-tuning run produces and consumes W\&B artifacts following
the same lineage pattern as the original pipeline. The Ascend
run is tagged with \texttt{["sft", "ascend", "telecom"]} in the
\texttt{patch\_wandb\_init\_for\_lineage()} call, enabling side-by-side
comparison of CUDA and Ascend runs in the W\&B dashboard. The
output artifact (\texttt{nano3-telecom-sft-model:latest}) is
consumed by both the inference compilation step and the loss
validation tool.

\begin{table}[htbp]
  \centering
  \caption{Artifact inputs and outputs for the Ascend SFT run}
  \label{tab:ascend-sft-artifacts}
  \begin{tabularx}{\textwidth}{llX}
    \toprule
    \tableheader{Direction & Artifact & Description} \\
    \midrule
    Input (model)
      & \texttt{nano3-pretrain-model:latest}
      & Pretrained Nemotron 3 Nano base checkpoint \\
    \rowcolor{TableRowOdd}
    Input (data)
      & \texttt{nano3-sft-data-telecom:latest}
      & Packed Parquet shards from telecom SFT data prep \\
    Output
      & \texttt{nano3-telecom-sft-model:latest}
      & Fine-tuned Megatron checkpoint; saved every 200 iterations \\
    \bottomrule
  \end{tabularx}
\end{table}
